.bp odd
.pn 1
.cs 5
.mt 5
.mb 6
.pl 66
.ll 65
.po 10
.hm 2
.fm 2
.he
.ft                                2-%
.ce 
.sh
Section 2
.qs
.sp
.ce
.sh 
Program Structure 
.qs
.tc 2  Program Structure 
.in 0
.sp 3
.sh
2.1  High-level Organization     
.qs
.ix high-level organization of PL/I programs
.he PL/I Reference Manual                2.1  High-level Organization        
.tc    2.1  High-level Organization
.pp 0
The following statements comprise every PL/I program:
.sp
.nf
.in 3
o Structural statements 
o Declarative statements 
o Executable statements 
.ix structural statements
.ix statements, structural
.fi
.in 0
.pp
Structural statements define distinct, logical units within a program and 
therefore determine the overall, high-level organization.  When a program
runs, control always flows from one of these logical units to another.  
Logical units can contain other logical units; they can be nested. 
Structural statements also determine the hierarchical structure of a program
where some logical units are subordinate to others.
.ix logical units
.ix declarative statements
.ix statements, declarative
.ix environment
.ix context for executable statements
.pp
Declarative statements determine the environment of a logical unit.  The 
environment is simply the names and attributes of variables that are available 
or active in a logical unit.  Declarative statements specify the context of 
variables that can be legally manipulated in a logical unit.
.pp
.ix executable statements
.ix statements, executable
Executable statements are statements that perform some action.  Both structural
statements and declarative statements serve only to create a context for 
executable statements.  All executable statements fall into one of the 
following categories:
.sp 2
.in 5
.ix preprocessor statements
.ix statements, preprocessor
.ti -2
.ix assignment statements
.ix statements, assignment
o Assignment statements that assign the value of an expression or constant to
a variable. 
.sp
.ti -2
.ix condition handling statements
.ix statements, condition handling
.ix run-time errors, interception of and recovery from
.ix error recovery
o Condition handling statements that allow a program to intercept and recover 
from run-time errors.
.sp
.ti -2
.ix I/O statements
.ix statements, I/O
o I/O statements that control the flow of data to and from I/O devices.
.sp
.ti -2
.ix memory management statements
.ix statements, memory management
o Memory management statements that manipulate storage.
.sp
.ti -2
.ix null statements
.ix statements, null
o Null statements that perform no action but function as placeholders. 
.sp
.ti -2
o Preprocessor statements that execute at compile time and manipulate external 
source files.
.sp
.ti -2
.ix sequence control statements
.ix statements, sequence control
.ix flow of control between logical units
o Sequence control statements that transfer the flow of control between 
logical units.
.sp
.in 0
.pp
.ix statement label
.ix PL/I keyword
All PL/I statements, except the assignment statement, consist of an optional
label, followed by a keyword (see section 2.2) and statement body, and end 
with a semicolon.  Subsequent sections of this manual describe each type of 
statement in detail.  For reference, Section 15 contains a complete 
alphabetical list of PL/I statements. 
.sp 2
.cp 5
.sh
2.2  Blocks
.he PL/I Reference Manual                                 2.2  Blocks 
.tc    2.2  Blocks
.ix block-structure in PL/I
.ix logical units
.pp
PL/I is a block-structured language.  That is, you group one or more
statements into logical units called blocks.  A block is a collection of 
statements in which declared variables are known.  Inside a block, you can 
declare variables and, for certain variables, you can allocate and fre� 
storage�  Yo� ca� nes� block� i� on� another� bu� no� overlap them.
.pp
.ix BEGIN blocks
.ix PROCEDURE blocks
.ix statements, BEGIN
.ix BEGIN statements
.ix statements, END
.ix END statements
.ix PROCEDURE statements
.ix statements, PROCEDURE
There are two types of blocks: BEGIN blocks and PROCEDURE blocks.  A BEGIN 
block is a sequence of statements delimited by BEGIN and END statements.  A 
PROCEDURE block is delimited by PROCEDURE and END statements.
.pp
A BEGIN block has the format:
.sp 2
.nf
.in 5
[label:]            
       BEGIN;       
       statement-1; 
            .       
            .       
            .       
       statement-n; 
       END [label]; 
.fi
.in 0
.sp 2
.ix nested BEGIN blocks
.ix block balance
where statement-1 through statement-n are any PL/I statements constituting 
the body of the block.  BEGIN blocks can contain nested PROCEDURE blocks, and 
nested BEGIN blocks.  I� PL/I� th� labe� optio� fo� th� END statement doe� no� 
automaticall� cause the bloc� t� balance� as it does in some implementations
of full PL/I.
.sp
.sh
Note:  \c
.qs
the label for the END statement is optional, but if included it must match the
label for the BEGIN statement.
.pp
A PROCEDURE block has the format:
.sp 2 
.nf
.in 5
proc-name:
         PROCEDURE-statement;
         statement-1;
              . 
              . 
              .
         statement-n;
         END [proc-name];
.fi
.in 0
.bp
where proc-name identifies the procedure, and statement-1 through statement-n 
are any PL/I statements constituting the body of the block.  Section 2.9
describes the PROCEDURE statement.  
.sp
.sh
Note:\c
.qs
\  the proc-name for the END statement is optional, but if included it must 
match the proc-name for the PROCEDURE statement.
.ix control
.pp
The essential difference between a BEGIN block and a PROCEDURE block is how
they receive control when the program is running.  Control flows into a BEGIN 
block in the usual sequential manner.  At this point, the block becomes 
active.  When control transfers, programmatically, outside the block, or its 
corresponding END statement executes, the block terminates.
.pp
.ix block activation
.ix block termination
.ix procedure invocation
.ix invoking a procedure
PL/I skips PROCEDURE blocks during the usual execution sequence, and they 
receive control only when invoked (see Section 2.5).  Figure 2-1 illustrates 
the block concept. 
.bp
.in 3
.nf
         PROCEDURE block                     BEGIN block

  A:                                      begin;
     procedure options(main);                     .
              .                                  .
              .                                  .
          statement(s)                         statement(s)
              .                                  .
              .                                  .
              .                                  .
         
  end A;                                  end;




        A:
          procedure options(main); 
              .
              .
              .
 
              begin;
               .
               .
               .
              end;


              begin; 
               .
               .
               .
              

                       begin;
                        .
                        .
                        .
                       end;

             end;

       end A;
.fi
.in 0
.sp 2
.ce
.sh
Figure 2-1.  BEGIN and PROCEDURE Blocks
.ix internal block
.ix external block
.ix nested block
.ix encompassing block
.bp 
.sh
2.3  Internal vs. External Blocks    
.qs
.he PL/I Reference Manual           2.3  Internal vs. External Blocks    
.tc    2.3  Internal vs. External Blocks
.ix main procedure
.pp
Each block is characterized as either internal or external depending on 
its relationship with other blocks.  An internal procedure is one that is 
contained in an encompassing block.  An external procedure is separate from 
other blocks.  The procedure is not contained (nested) in any other block.  
Thus, the main procedure is always an external procedure.  
.pp
A PL/I program can have one or more external procedures that contain nested 
internal procedures or blocks.  Each external procedure can be separately 
compiled and linked together to form a runnable program.  One of the external 
procedures forming the program must be the main procedure.  
.ix separate compilation of external procedures
.pp
In Figure 2-2(a), blocks P1, P2, and P3 are all external but the BEGIN block 
is internal to P3.  In Figure 2-2(b), P1 is the external block, and P2, P3, 
and the BEGIN block are all internal.  The main procedure has the form:
.sp 2
.in 5
.nf
proc-name:
      PROCEDURE OPTIONS(MAIN);
               .
               .
               .
      Statements or Blocks
               .
               .
               .
      END [proc-name];
.in 0
.bp
.nf

    P1:                             P1:
      procedure options(main);        procedure options(main);
             
             
             
    end P1;                                   


    P2:                                 P2:           
      procedure;                          procedure;
            
            
    end P2;                             end P2;


    P3:                                P3:
      procedure;                         procedure;
            
            

          begin;                             begin;
            
            
            
          end;                               end;

          
          
                                       end P3;
    end P3;                                                       
                                    end P1;                      

.sh
               (a)                               (b)

.fi
.sp 2
.ce
.sh
Figure 2-2.  Internal and External Blocks
.sp 2
.ix external block
.ix internal block
.pp
The 
.ul
PL/I Language Programmer's Guide\c 
.qu
\ contains specific examples of program structure and how you can separately 
compile, link, and load external procedures.
.ix external procedures
.sp 2
.sh
2.4  Scope of Variables 
.qs
.he PL/I Reference Manual                     2.4  Scope of Variables 
.tc    2.4  Scope of Variables 
.ix scope of a variable
.ix variables, local
.ix varibles, external
.ix containing block
.ix local variable
.ix external variable
.pp
The scop� o� � variabl� i� th� set of blocks in which the variable is known.
Variables can be either local or external relative to a block in which they 
appear.  
.bp
.pp 0
When you declare a variable in a block, you can reference it in that block
or any contained block.  The variable is said to be local to that block 
because you cannot reference it outside the block where you declare it.  In a 
contained block, a reference to a variable declared in a containing block is 
called an up-level reference. 
.ix up-level reference
.pp
The following example illustrates the concept of scope:
.sp 2
.in 5
.nf
P1:
  procedure;                                                         
  declare                                                            
  (a,b) fixed binary(7);                                             
  a = 2;       /* a is local to P1 */                                
  b = 3;       /* b is local to P1 */                                
  P2:                                                                
    procedure;                                                       
    declare                                                          
    b fixed binary(7);                                               
    b = 2;   /* b is local to P2 */                                  
    a = a*b; /* b here is b in P2, not b in P1 */                    
  end P2;                                                            
  put list (a,b);                                                    
                                                             
end P1;                                                            
.fi
.in 0
.sp  
.pp
PL/I creates a new variable b in block P2 because it is a declared variable in 
that block.  The PUT LIST statement is outside P2; therefore, the value of the 
variable b of P1 is 3.  Because there is no declaration for the identifier a 
in P2, a is an up-level reference to the variable a declared in P1, and the 
assignment statement in P2 changes its value.  Thus, this code sequence 
produces the values 4 and 3.
.pp
Any variable declared as EXTERNAL is known to all blocks in which it is 
declared as EXTERNAL and in all contained blocks unless redeclared without the 
EXTERNAL attribute.  Two declarations of the same variable name denote 
separate storage locations unless both specify the EXTERNAL attribute.
.ix EXTERNAL attribute
.bp
.in 5
.nf
P1:                                                            
   procedure;                                                  
   declare                                                     
   z fixed binary external;                                    
            .                                                  
            .                                                  
            .                                                  
   P2:                                                         
      procedure;                                               
      declare                                                  
      z fixed binary external;                                 
               .                                               
               .                                               
               .                                               
      P3:                                                      
         begin;                                                
         declare                                               
         z float binary; /* not external */                    
                  .                                            
                  .                                            
         end;                                                  
      end P3;                                                  
   end P2;                                                     
end P1;                                                                   
.fi
.in 0
.sp 
.pp
In this code sequence, the variable z in P1 and P2 refers to the same 
external variable, but variable z in P3 is a local variable and is distinct 
from the external variable z.  
.ix local variable
.ix external variable
.bp
.nf
.in 5
P1:                                                    
  procedure options(main);                             
   declare x float binary;                             
        .                                              
        .                                              
        .                                              
    begin;                                             
        declare x fixed;                               
            .                                          
            .                                          
            .                                          
            .                                          
    end;                                               
        .                                              
        .                                              
        .                                              
  P2:                                                  
    procedure;                                         
    declare x character(10) varying;                   
         .                                             
         .                                             
         .                                             
  end P2;                                              
                                               
end P1;                                                
.sp 
.in 0
.fi
.pp
In this code sequence, the scope of x is limited to each block in which it is 
declared.  Although the name is identical in each declaration, the compiler
treats each one as a completely different variable with its own data type, and
stores them in different memory locations.
.ix scope of a variable
.ix compiler
.sp 2
.tc    2.5  Procedure Blocks
.ix procedure blocks
.ix subroutine
.ix argument-list, for a procedure 
.ix procedure invocation
.ix invoking a procedure
.sh
2.5  Procedure Blocks
.qs
.he PL/I Reference Manual                       2.5  Procedure Blocks
.pp 0
In PL/I, there are two types of procedures:  subroutines and functions.  Both 
types perform a specific task and are logically separate from the rest of the 
program.  Both types can execute the same sequence of code one or more times 
without duplicating the code at each occurrence.  
.pp
You invoke or call a subroutine and, optionally, pass data items to it in an 
argument list.  The subroutine then manipulates the data and, optionally, 
returns it to the invoking procedure.  Control resumes at the statement
immediately following the invocation. 
.ix function 
.ix flow of control within a program
.ix argument-list, for a function
.pp
A function is a procedure that manipulates data items and then returns a 
single value.  You invoke a function by referencing its function name and 
argument-list in an expression.  Control passes to the function that performs 
its task and then returns a single value that replaces the function 
reference.  Control then resumes at the point of reference.  
.bp
.sh
2.6  The CALL Statement
.qs
.tc    2.6  The CALL Statement
.he PL/I Reference Manual                     2.6  The CALL Statement
.ix CALL statement
.pp
The CALL statement has the general form:
.sp
.ti 5
CALL proc-name[(sub-1,...,sub-n)] [(argument-list)];
.sp
.mb 4
.fm 1
where sub-1 through sub-n are optional subscripts that are required only when
proc-name is a subscripted entry variable (Section 3.3.2), and argument-list 
represents the arguments passed to the procedure.  Figure 2-3 
illustrates the invocation of subroutines and functions.
.ix subscripted entry variable
.ix arguments
.ix argument-list, in a CALL statement
.ix procedure invocation
.ix subroutine invocation
.ix invoking a procedure
.ix invoking a subroutine
.sp 3
.nf
.cp 14
Subroutine invocation:               Function invocation:
 
CALL name [argument-list];           name [argument-list] 


Subroutine name   arguments          function name   arguments
.sp 2
.ce
.sh
Figure 2-3.  Subroutine and Function Invocation
.sp 3
example:                             example:

call print_header;                   point = 3.14/sin(A);
call compute(base_pay,overtime);     put list (sum(X,Y));
.fi
.sp 2
.tc    2.7  The RETURN Statement
.sh
2.7  The RETURN Statement
.qs
.he PL/I Reference Manual                   2.7  The RETURN Statement
.pp
.ix RETURN statement
.ix statements, RETURN
The RETURN statement returns control to the point in the calling block 
immediately following the procedure invocation.  It also returns a value if 
the procedure is a function procedure.
.pp
.ix procedure invocation
.ix function procedure
The RETURN statement has the form:
.sp 
.ti 5
RETURN [(return-exp)];
.sp 
where return-exp is the function value the procedure returns to the calling 
point.  When necessary, PL/I converts the attributes of the returned value to 
conform to the attributes specified in the RETURNS attribute of the procedure 
statement.  (See Section 4.1.)
.pp
The RETURN statement ends the procedure block that contains it.  If the main 
procedure has the RETURNS attribute, PL/I returns control to the operating 
system.
.pp
.ix operating system
.ix procedure block
.ix RETURN statement
.ix statements, RETURN
.ix RETURNS attribute
The following are some examples of RETURN statements: 
.sp
.nf
.in 5
return;
return (X**2);
return (F(A,(B)));
.fi
.in 0
.bp
.tc    2.8  Arguments and Parameters
.he PL/I Reference Manual           2.8  Arguments and Parameters
.sh
2.8  Arguments and Parameters
.qs
.ix arguments, values passed to a procedure
.ix parameters, values expected by a procedure
.pp
The data items you pass to a procedure are called the arguments,
while the data items expected by a procedure and defined in the 
PROCEDURE statement, are called the parameters.  Upon invocation of a 
procedure block, PL/I pairs each argument with its corresponding 
parameter.  Figure 2-4 illustrates this concept.
.ix arguments
.ix PROCEDURE statement
.ix statements, PROCEDURE
.ix parameters
.mb 6
.fm 3
.sp 2
.nf
                                       argument list
    call compute (a+(b+c),r/2,3.14);
               .
               .

    compute: procedure (X,Y,Z);
               .
               .                       parameter-list
               .
    end compute;        
.sp 2
.ce
.sh
Figure 2-4.  Arguments and Parameters
.qs
.sp 2
.fi
.pp
When you pass the argument by reference, the argument and 
corresponding parameter share storage.  In this case, any changes made 
to the parameter in the invoked procedure change the value of argument of the 
invoking block.  
.ix parameter passing, by reference
.ix parameter passing, by value
.pp
When you pass the argument by value, the argument and parameter do not share 
storage.  In this case, PL/I passes a copy of the argument to the invoked 
procedure, so that any changes to the 
parameter affect only the copy, not the argument's value.
.ix procedure invocation
.ix invoking a procedure
.ix storage sharing
.mb 4
.fm 1
.pp
The following example program illustrates parameter passing.
.sp 2
.nf
.in 5
A:                               
 procedure;                              
 declare                                 
    ACTUAL fixed binary,                 
    DUMMY  fixed binary;                 
                                         
 call X(ACTUAL);                 
 call X((DUMMY));                        
                                         
 X:                              
  procedure (FORMAL);                    
  declare FORMAL fixed binary;           
  FORMAL = 3;                            
 end X;                                  
                                         
end A;                            
.in 0        
.bp
.fi
.pp
PL/I passes ACTUAL by reference.  Therefore, the assignment statement in 
the procedure X changes the value of ACTUAL throughout the program.  PL/I
passes DUMMY by value.  Thus the procedure only changes a copy of the value  
inside the procedure.
.pp
PL/I passes arguments by reference when the data attributes of the 
argument are the same as the data attributes of the parameter.  
PL/I passes an argument by value when it is one of the following:
.ix data type matching when passing parameters
.sp 2
.in 5
.ti -2
o a constant 
.ti -2
o an entry name
.ti -2
o an expression consisting of variable references and operators
.ti -2
o a variable reference enclosed in parentheses 
.ti -2
o a function invocation
.ti -2
o a variable reference whose data type does not match that of the parameter
.sp 2
.in 0
In the latter case, PL/I converts the argument to the data type, precision, 
and scale factor of the parameter.  The following program illustrates 
this concept:
.ix precision
.ix scale factor
.mb 4
.fm 1
.nf
.sp 2
.in 5
A:                                   
 procedure;                          
 declare                             
   X character(7),                  
   (Y,Z) fixed binary;               
                             
 call p(X,(Y),Z);                    
                             
 p:                                  
     procedure(A,B,C);               
     declare                         
       A character(7),              
       B fixed binary,               
       C float binary;               
                                    
     A = 'Digital';                  
     B = 100;                        
     C = 2.5E2;                      
                                    
 end p;                              
                             
end A;                               
.in 0
.sp 
.fi
.pp
The CALL statement sends the procedure three arguments X, Y, and
Z corresponding to the three parameters A, B, and C.  PL/I passes the 
first argument by reference because it matches the parameter, 
and the second argument by value because it occurs as an expression.  PL/I 
converts the third argument to the FLOAT binary data type and passes it by 
value.
.ix CALL statement
.ix statements, CALL
.bp
.tc    2.9  The PROCEDURE Statement
.ix PROCEDURE statement
.ix statements, PROCEDURE
.ix readability of programs
.ix main program
.ix main procedure statement
.he PL/I Reference Manual                2.9  The PROCEDURE Statement
.sh
2.9  The PROCEDURE Statement
.qs
.pp
In PL/I, you can define a procedure with a PROCEDURE statement at any point in 
a program.  However, for readability you should place all procedures together 
in a single section at the beginning or the end of the main program.  The 
main program is a single-procedure definition.  
.pp
The PROCEDURE statement identifies the entry point to the procedure, delimits 
the beginning of the procedure block, defines the parameter list, and 
gives the attributes of the returned value for functions.  The procedure can 
consist of a sequence of one or more statements including the corresponding 
END statement that ends the procedure definition.  The END statement can also 
be the exit point of the procedure, although embedded RETURN statements can 
appear within the procedure body.
.ix END statement
.ix statements, END
.ix statements, RETURN
.ix RETURN statement
.ix procedure entry point
.ix procedure exit point
.pp
The PROCEDURE statement has the general form:
.sp
.mb 6
.fm 3
.in 5 
.nf
proc-name: PROCEDURE[(parameter-list)]
           [OPTIONS(option,...)] [RETURNS(attribute-list)] 
           [RECURSIVE];
.fi
.in 0
.sp
where parameter-list are the parameters for the procedure which
you must declare within the procedure body at the principle block level.  A 
parameter can be any of the following:
.ix parameter
.ix parameter-list, in a PROCEDURE statement
.sp
.nf
.in 3
o a scalar variable
o an array 
o a major structure 
.fi
.in 0
.sp 
but cannot have the attributes:
.sp
.nf
.in 3
o STATIC   
o AUTOMATIC   
o BASED   
o EXTERNAL
.fi
.in 0
.pp
OPTIONS(option,...) defines a list of one or more of the options MAIN,
STACK(b), or EXTERNAL.  
.sp 2
.in 5
.ti -2
o The MAIN option identifies the procedure as the first procedure to receive 
control when the program begins execution. 
.ix MAIN option, in a PROCEDURE statement
.sp
.ti -2
o The STACK(b) option sets the size of the run-time stack to the number of bytes 
specified by b.  The default value is 512 bytes.  
.ix STACK(b) option, in a PROCEDURE statement
.ix run-time stack, default value
.ix default value, of run-time stack
.sp
.ti -2
o The EXTERNAL option identifies the procedure as an externally compiled
procedure.  The EXTERNAL option in a procedure heading makes the procedure 
accessible outside the module.  It is often useful to group separately 
compiled procedures into a 
single compilation, where the procedures reference the same global data.  
According to the Subset G standard, you must compile each subroutine 
separately, and duplicate the global data area in each compilation.  You can
then combine the individual modules using the linkage editor to produce the 
object module.
.in 0
.ix EXTERNAL option, in a PROCEDURE statement
.ix global data in PL/I programs
.ix Subset G standard
.sp
.pp
The following procedure shows an example of using the EXTERNAL option:
.sp 2
.nf
.in 5
module:                                                      
    procedure;                                               
    declare                                                  
       1 global_data static,                                 
         2 a_field character(20) varying initial(''),        
         2 b_field fixed initial(0),                         
         2 c_field float initial(0);                         
    set_a:                                                   
      procedure (c) options(external);                       
      declare c character(20) varying;                       
      a_field = c;                                           
    end set_a;                                               
    set_b:                                                   
      procedure (x) options(external);                       
      declare x fixed;                                       
      b_field = x;                                           
    end set_b;                                               
    set_c:                                                   
      procedure (y) options(external);                       
      declare y float;                                       
      c_field = y;                                           
    end set_c;                                               
    sum:                                                     
      procedure returns(float) options(external);            
      return (b_field + c_field);                            
    end sum;                                                 
    display:                                                 
      procedure options(external);                           
      put skip list(a_field,b_field,c_field);                
    end display;                                             
end module;                                                  
.in 0
.bp
.fi
.pp
This code defines five external procedures: set_a, set_b, set_c, sum, and
display.  These procedures are then accessed in the following code sequence:
.sp 2
.in 5
.nf
call_ext:                                          
   procedure options(main);                        
   declare                                         
      set_a entry (character(20) varying),         
      set_b entry (fixed),                         
      set_c entry (float),                         
      sum returns(float),                          
      display entry;                               
   call set_a('Johnson,J');                        
   call set_b(25);                                 
   call set_c(5.50);                               
   put skip list(sum());                           
   call display();                                 
end call_ext;                                      
.fi
.in 0
.sp 
.pp
These two modules, when compiled separately and linked together, form a 
single, runnable program.
.sp 2
.in 5
.ti -2
o The RETURNS attribute for a function procedure gives the attributes of the 
value returned by the function.
.ix RETURNS attribute
.sp
.ti -2
o The RECURSIVE attribute indicates that the procedure can activate itself, 
either directly or indirectly.
.in 0
.ix RECURSIVE attribute
.nx twob.tex


